\documentclass[11pt]{article}
\usepackage{fontspec,amsmath,amssymb,amsthm,hyperref}
\usepackage[margin=1in]{geometry}
\setlength{\emergencystretch}{3em}
\title{Contraction of the actual KIPLMC2 steps and chains}
\author{LangevinLean mathematical exposition}
\date{}
\begin{document}\maketitle
\noindent Source module: \texttt{KIPLMC/Contraction.lean}. Compile this proof document with XeLaTeX or LuaLaTeX. The accompanying HTML pairs the mathematical argument with the exact Lean statements and records the scope of the formalization.


KIPLMC2 inherits the shared OBABO contraction by using the actual averaged particle potential. Its gradient is exactly the implemented particle force, and its strong convexity and smoothness constants remain \(m,L\), independent of the number of particles. The thermostat temperature changes the noise amplitude; common innovations cancel exactly in a synchronous comparison.


\section{Transfer the proved one-step phase inequality}


Write \(\mathcal E(v,x)=\lVert x\rVert^2+2\langle x,v\rangle/\gamma+\lVert v\rVert^2/L\). Under \(\gamma^2\ge4L\) and \(33\gamma^3h\le m\), the shared theorem proves \(\mathcal E(S(z,\xi,\xi\prime)-S(w,\xi,\xi\prime))\le(1-mh/(3\gamma))\mathcal E(z-w)\). Lean supplies the proved differentiability, strong convexity and Lipschitz gradient of the average potential, then rewrites its gradient as the actual force. Both thermostat inputs are the same in the two copies. No Gaussian assumption is needed for this pathwise inequality.



\noindent Lean declarations: \texttt{kiplmc2Step\_synchronous\_contraction}.

\section{Use the implemented chain recurrence at every step}


The chain is the concrete noise-driven recursion, not a sequence assumed to satisfy an estimate. Its successor identity supplies exactly the shared OBABO iteration hypotheses. Repeated contraction gives \(\mathcal E(Z_n-\widetilde Z_n)\le q^n\mathcal E(Z_0-\widetilde Z_0)\), where \(q=1-mh/(3\gamma)\). All noise values may be arbitrary: the result therefore applies to the Gaussian algorithm after its independent-noise law is constructed.



\noindent Lean declarations: \texttt{kiplmc2Chain\_synchronous\_contraction}.

\section{Project to positions with a common initial velocity}


The phase energy satisfies \(\mathcal E(v,x)\ge\lVert x\rVert^2/2\). If the initial velocities agree, its initial value is exactly the squared initial position discrepancy. Taking square roots gives \(\lVert X_n-\widetilde X_n\rVert\le\sqrt{2q^n}\lVert X_0-\widetilde X_0\rVert\). The physical-time contraction does not acquire an artificial particle factor. This is the stability part of the sampling analysis; comparison with the Gibbs target also needs the separately proved discretization bias.



\noindent Lean declarations: \texttt{kiplmc2Chain\_position\_decay}.

\end{document}
